\section{Limitations}

\paragraph{Single seed.} One run per condition. An observed gap could be seed noise rather
than a property of the mechanism. Multi-seed replication is the first thing to add.

\paragraph{Single scale.} One dataset, one model size, one 32{,}000-step budget. A null
result does not disprove the effect at larger scale; it bounds where we looked. The
secondhand result being replicated was reported around 20{,}000 steps at an unknown model
scale, so our budget may sit near or below the threshold where the effect appears.

\paragraph{Unvalidated SIGReg implementation.} \interp{} Written from a summary-level
understanding of \citet{lejepa2025}, not validated against the reference implementation. A
negative result could reflect an implementation discrepancy rather than a property of
SIGReg.

\paragraph{Simplified predictor.} An MLP conditioned on target position, rather than the
narrow transformer used by \citet{assran2023ijepa}.

\paragraph{Departure detection is crude.} The departure step uses a fixed threshold against
an early-training baseline. It is applied identically to every series, which is what the
comparison requires, but it is not a calibrated changepoint test.

\paragraph{The instrument finding rests on one pilot.} \interp{} The scale-invariance
failure reported in Section~\ref{sec:instrument} was observed in a single 2{,}000-step
run at one model scale, on one seed. The mechanism is straightforward and we reproduced
the probe half of it synthetically, but the magnitude of the effect at other scales,
architectures, and embedding dimensionalities is unmeasured. We also have not surveyed
the literature systematically enough to claim the joint blind spot is unreported.

\paragraph{Pilot orderings are not results.} No between-condition comparison from the
pilot should be read as evidence about the research question. Those numbers were produced
by the instrument that Section~\ref{sec:instrument} shows was defective.

\paragraph{Everything before Phase 2 is single-seed and no-projector.} The Part-1 and
Part-2 results above await the pre-registered multi-seed rerun with the completed condition
grid; until then they are supporting evidence for the instrument findings, not standalone
claims.
